\chapterbackmatter{Solutions to Supplementary Exercises}

\begin{enumerate}
\SupplementarySolution{1}{M{\o}ller Operators and Scattering States}
For fixed \(s\), shifting the limiting variable gives
\[
 e^{iHs}\Omega_\pm e^{-iH_0s}
 =\lim_{t\to\mp\infty}e^{iH(t+s)}e^{-iH_0(t+s)}
 =\Omega_\pm .
\]
Differentiation at \(s=0\) yields the intertwining relation
\[
 H\Omega_\pm=\Omega_\pm H_0.
\]
Thus an \(H_0\)-eigenstate of energy \(E_\alpha\) is mapped to a
generalized \(H\)-eigenstate of the same energy.  With the signs in the
question,
\[
 \InKet{\alpha}=\Omega_+\ket{\alpha}_0,\qquad
 \OutKet{\alpha}=\Omega_-\ket{\alpha}_0.
\]
On the scattering subspace the wave operators are isometries.  The
operator acting on free asymptotic labels is
\[
 S=\Omega_-^\dagger\Omega_+,\qquad
 \braket{\beta}{S\alpha}_0=\OutBra{\beta}\InKet{\alpha}.
\]
If their ranges exhaust the continuous spectral subspace, \(S\) is
unitary there; bound states lie outside this scattering subspace.

\SupplementarySolution{2}{The \(S\)-Matrix as an In--Out Overlap}
The total momentum acts on the asymptotic states with eigenvalues
\(p_\alpha\) and \(p_\beta\).  Translation by \(a\) therefore gives
\[
 \OutBra{\beta}\InKet{\alpha}
 =e^{i(p_\beta-p_\alpha)\cdot a}
 \OutBra{\beta}\InKet{\alpha}
\]
for every \(a\).  The overlap consequently has support only at
\(p_\beta=p_\alpha\).  Isolating free propagation in the book's
normalization gives
\[
 \OutBra{\beta}\InKet{\alpha}
 =\delta(\beta-\alpha)
 -2\pi i\,\delta^4(p_\beta-p_\alpha)M_{\beta\alpha}.
\]
The identity term contains delta functions matching every outgoing
particle with an incoming particle.  A connected matrix element retains
only the single delta function required by conservation of the total
four-momentum.  Extra delta functions matching proper subsets signal
disconnected spectator pieces.

\SupplementarySolution{3}{Completeness on the Scattering Subspace}
Symmetrized two-boson states obey
\[
 \braket{\mathbf k,\mathbf l}{\mathbf p,\mathbf q}
 =\delta^3(\mathbf k-\mathbf p)\delta^3(\mathbf l-\mathbf q)
 +\delta^3(\mathbf k-\mathbf q)\delta^3(\mathbf l-\mathbf p).
\]
Hence
\[
 I_2=\frac12\int d^3p\,d^3q\,
 \ket{\mathbf p,\mathbf q}\bra{\mathbf p,\mathbf q}
\]
acts on \(\ket{\mathbf k,\mathbf l}\) as one half times two equal terms,
and is the identity.  Equivalently one may integrate over only one
ordering of the two momenta and omit \(1/2!\).

Applying the wave operators gives
\[
 I_{2,\rm in}=\frac12\int d^3p\,d^3q\,
 \InKet{\mathbf p,\mathbf q}\InBra{\mathbf p,\mathbf q},
\quad
 I_{2,\rm out}=\frac12\int d^3p\,d^3q\,
 \OutKet{\mathbf p,\mathbf q}\OutBra{\mathbf p,\mathbf q}.
\]
These are identities on the corresponding two-particle asymptotic
subspaces.  Asymptotic completeness asserts that the sum over every
open particle-number channel supplies the full scattering projector.

\SupplementarySolution{4}{Invariant Relative Velocity}
With \(p_i^2=-m_i^2\), the numerator
\[
 F=\sqrt{(p_1\cdot p_2)^2-m_1^2m_2^2}
\]
is invariant.  The ratio \(F/(E_1E_2)\) is the relative velocity that
combines with the density factors implicit in the book's
\(\delta^3\)-normalized states.  If \(p_2=(m_2,\mathbf0)\), then
\[
 F=m_2|\mathbf p_1|,\qquad
 u=\frac{|\mathbf p_1|}{E_1}=|\mathbf v_1|.
\]
For collinear velocities,
\[
 \frac{F^2}{E_1^2E_2^2}
 =(1-v_1v_2)^2-(1-v_1^2)(1-v_2^2)
 =(v_1-v_2)^2.
\]
Thus \(u=|v_1-v_2|\), including the case of oppositely directed
particles through the signs of \(v_1,v_2\).

\SupplementarySolution{5}{Two-Body Phase Space in the Center-of-Mass Frame}
Momentum conservation sets
\(\mathbf p'_2=-\mathbf p'_1\).  Write \(q=|\mathbf p'_1|\).  The
remaining energy delta function is
\[
 \delta(E-E'_1(q)-E'_2(q)),
\quad
 \frac{d(E'_1+E'_2)}{dq}
 =\frac{qE}{E'_1E'_2}.
\]
Consequently
\[
 \delta^4(p_\beta-p_\alpha)d\beta
 \longrightarrow \frac{k'E'_1E'_2}{E}\,d\Omega,
\]
which is Eq.~\eqref{eq:3.4.28}.  Squaring
\(E=E'_1+E'_2\) and solving for the positive momentum gives
\[
 k'=\frac{\sqrt{[E^2-(m'_1+m'_2)^2]
 [E^2-(m'_1-m'_2)^2]}}{2E}.
\]
The expression vanishes at threshold and is real precisely when the
final channel is open.

\SupplementarySolution{6}{A General Two-Body Decay Width}
Equation \eqref{eq:3.4.29} gives
\[
 \frac{d\Gamma}{d\Omega}
 =\frac{2\pi k'E'_1E'_2}{M}|M_{fi}|^2.
\]
For an angle-independent matrix element,
\[
 \Gamma
 =\frac{8\pi^2k'E'_1E'_2}{M}|M_{fi}|^2,
\]
where
\[
 k'=\frac{\sqrt{[M^2-(m_1+m_2)^2]
 [M^2-(m_1-m_2)^2]}}{2M},
\quad
 E'_{1,2}=\frac{M^2+m_{1,2}^2-m_{2,1}^2}{2M}.
\]
These factors are specific to the book's \(\delta^3\) state normalization.
For two identical daughters, integrating over the full labeled phase
space double-counts every physical state, so the result is divided by
\(2!\).

\SupplementarySolution{7}{Optical Theorem from \(S^\dagger S=1\)}
Substitution of \(S=1-i\mathcal T\) into unitarity yields
\[
 i(\mathcal T-\mathcal T^\dagger)=\mathcal T^\dagger\mathcal T.
\]
Between arbitrary states,
\[
 i(\mathcal T_{fi}-\mathcal T_{if}^*)
 =\sum_n\mathcal T_{nf}^*\mathcal T_{ni},
\]
where the sum includes all discrete labels and phase-space integrals.
For \(f=i\),
\[
 2\,\operatorname{Im}\mathcal T_{ii}
 =\sum_n|\mathcal T_{ni}|^2.
\]
Every term on the right is a squared modulus, so the absorptive part of
the forward matrix element is nonnegative.  This is the optical theorem
before the common energy--momentum delta functions are removed.

\SupplementarySolution{8}{Forward Scattering and the Total Cross Section}
In the book's convention,
\[
 \mathcal T_{\beta\alpha}
 =2\pi\delta^4(p_\beta-p_\alpha)M_{\beta\alpha}.
\]
Removing the common factor
\(\delta^4(0)\) from the forward optical theorem gives
\[
 \operatorname{Im}M_{\alpha\alpha}
 =\pi\int d\beta\,
 \delta^4(p_\beta-p_\alpha)|M_{\beta\alpha}|^2.
\]
On the other hand, integration of Eq.~\eqref{eq:3.4.15} gives
\[
 \sigma_{\rm tot}(\alpha)
 =\frac{(2\pi)^4}{u_\alpha}
 \int d\beta\,\delta^4(p_\beta-p_\alpha)
 |M_{\beta\alpha}|^2.
\]
Therefore
\[
 \boxed{\;
 \sigma_{\rm tot}(\alpha)
 =\frac{16\pi^3}{u_\alpha}
 \operatorname{Im}M_{\alpha\alpha}.
 \;}
\]
The numerical factor changes when covariantly normalized one-particle
states and the more common invariant amplitude are used; the content of
the theorem does not.

\SupplementarySolution{9}{Exact Scattering from a Delta Potential}
For incidence from the left, take
\[
 \psi_L(x)=
 \begin{cases}
 e^{ikx}+r e^{-ikx},&x<0,\\
 t e^{ikx},&x>0.
 \end{cases}
\]
Continuity gives \(t=1+r\), while
\(\psi'(0^+)-\psi'(0^-)=2mV_0\psi(0)\) gives, with
\(\lambda=mV_0/k\),
\[
 t=\frac1{1+i\lambda},\qquad
 r=\frac{-i\lambda}{1+i\lambda}.
\]
Reflection symmetry gives the right-incident solution by \(x\mapsto-x\).
In the basis of incoming waves from the left and right,
\[
 S=\begin{pmatrix}t&r\\ r&t\end{pmatrix}.
\]
Direct calculation gives
\[
 |t|^2+|r|^2=1,\qquad tr^*+rt^*=0,
\]
so \(S^\dagger S=1\).  As \(V_0\to0\), \(t\to1\); as
\(V_0\to\infty\), \(r\to-1\), as expected.

\SupplementarySolution{10}{Disconnected Spectators}
Let \(s\) denote the spectator and \(\widehat\alpha,\widehat\beta\) the
interacting parts.  With the book's normalization, a disconnected term
has the form
\[
 \OutBra{\mathbf p'_s,s;\widehat\beta}
 \left.\InKet{\mathbf p_s,s;\widehat\alpha}\right|_{\rm disc}
 =\delta^3(\mathbf p'_s-\mathbf p_s)\,
 \OutBra{\widehat\beta}\InKet{\widehat\alpha},
\]
with an additional Kronecker delta for the spectator's spin and species.
The reduced overlap contains its own conservation delta function.  Thus
the full term is more singular than a connected amplitude, because a
proper subset of momenta is separately conserved.  Defining connected
matrix elements by cumulant subtraction removes every such factorization,
leaving only the one delta function for total four-momentum conservation.

\SupplementarySolution{11}{The Mandelstam Sum Rule}
Momentum conservation gives \(p_2=p_3+p_4-p_1\).  Expanding the
definitions,
\[
 \begin{split}
 s+t+u
 &=-3p_1^2-2p_1\cdot(p_2-p_3-p_4)-p_2^2\\
 &=-p_1^2-p_2^2-p_3^2-p_4^2,
 \end{split}
\]
where the second line follows after using conservation once more.
Because the mostly-plus mass shells obey \(p_i^2=-m_i^2\),
\[
 s+t+u=m_1^2+m_2^2+m_3^2+m_4^2.
\]
This proof does not require any particular Lorentz frame.

\SupplementarySolution{12}{Massless Center-of-Mass Kinematics}
Choose
\[
 p_1=\frac E2(1,0,0,1),\quad
 p_2=\frac E2(1,0,0,-1),\quad
 p_3=\frac E2(1,\sin\theta,0,\cos\theta),
\]
with \(p_4\) opposite to \(p_3\) and \(E=E_{\rm cm}\).  In the definitions
of S.3.11,
\[
 s=E^2,\qquad
 t=-\frac{E^2}{2}(1-\cos\theta),\qquad
 u=-\frac{E^2}{2}(1+\cos\theta).
\]
They sum to zero.  Since \(-1\leq\cos\theta\leq1\),
\[
 -s\leq t\leq0,\qquad -s\leq u\leq0,
\]
and \(t+u=-s\).

\SupplementarySolution{13}{Scalar Yukawa Decay}
The vertex gives the covariant invariant amplitude
\(\mathcal M=-g\), so \(|\mathcal M|^2=g^2\).  In the rest frame of
\(\phi\),
\[
 |\mathbf p|=\frac M2\sqrt{1-\frac{4m^2}{M^2}}.
\]
The standard two-body phase-space integral gives
\[
 \Gamma(\phi\to\psi\psi^*)
 =\frac{|\mathbf p|}{8\pi M^2}|\mathcal M|^2
 =\frac{g^2}{16\pi M}
 \sqrt{1-\frac{4m^2}{M^2}}.
\]
The daughters have opposite charge and are distinguishable, so there is
no \(1/2!\) factor.

\SupplementarySolution{14}{Two Channels of Scalar Yukawa Scattering}
Charge flow permits annihilation through the \(s\)-channel and exchange
through the \(t\)-channel.  In the usual covariant amplitude convention,
\[
 \mathcal M=-g^2\left(
 \frac1{s-M^2+i0}+\frac1{t-M^2+i0}\right).
\]
There is no \(u\)-channel graph for these assigned charges.  When
\(|s|,|t|\ll M^2\),
\[
 \mathcal M
 =\frac{2g^2}{M^2}
 +O\!\left(\frac{s+t}{M^4}\right).
\]
The leading term is momentum independent and hence can be reproduced by
a local four-field interaction.

\SupplementarySolution{15}{Integrating Out a Heavy Scalar}
Write \(X=\psi^*\psi\).  At momenta small compared with \(M\), the heavy
field equation is
\[
 (\Box+M^2)\phi=-gX,\qquad
 \phi=-\frac g{M^2}X+O(M^{-4}\Box X).
\]
Substitution into the mass and interaction terms gives
\[
 \mathcal L_{\rm eff}
 =\frac{g^2}{2M^2}X^2+O(M^{-4})
 =-\lambda_{\rm eff}(\psi^*\psi)^2+\cdots,
\qquad
 \lambda_{\rm eff}=-\frac{g^2}{2M^2}.
\]
Differentiating \(X^2\) with respect to two \(\psi\)'s and two
\(\psi^*\)'s supplies a factor \(2!\,2!=4\).  The contact vertex is
\(-4i\lambda_{\rm eff}=2ig^2/M^2\), exactly \(i\mathcal M\) from the
leading heavy-exchange result.

\SupplementarySolution{16}{An \(O(N)\)-Invariant Four-Scalar Vertex}
Four differentiation of
\(-(\lambda/4)(\Phi_a\Phi_a)^2\) gives the vertex
\[
 -2i\lambda\left(
 \delta_{ij}\delta_{kl}+\delta_{ik}\delta_{jl}
 +\delta_{il}\delta_{jk}\right).
\]
Thus, omitting the common \(i\) in \(i\mathcal M\),
\[
 \begin{array}{c|c}
 \text{process}&\mathcal M\\ \hline
 \Phi_1\Phi_2\to\Phi_1\Phi_2&-2\lambda\\
 \Phi_1\Phi_1\to\Phi_2\Phi_2&-2\lambda\\
 \Phi_1\Phi_1\to\Phi_1\Phi_1&-6\lambda
 \end{array}
\]
The enhancement in the last row is the sum of the three equal invariant
index contractions.

\SupplementarySolution{17}{Flavor and Identical-Particle Cross Sections}
For equal masses, constant \(\mathcal M\), and \(k'=k\),
\[
 \frac{d\sigma}{d\Omega}
 =\frac{|\mathcal M|^2}{64\pi^2s}.
\]
Integration over the full sphere gives
\(|\mathcal M|^2/(16\pi s)\) for distinguishable final particles and
half that for identical final particles.  Therefore
\[
 \begin{split}
 \sigma(\Phi_1\Phi_2\to\Phi_1\Phi_2)
 &=\frac{\lambda^2}{4\pi s},\\
 \sigma(\Phi_1\Phi_1\to\Phi_2\Phi_2)
 &=\frac{\lambda^2}{8\pi s},\\
 \sigma(\Phi_1\Phi_1\to\Phi_1\Phi_1)
 &=\frac{9\lambda^2}{8\pi s}.
 \end{split}
\]
There is no extra \(1/2\) for a specified incident pair; the statistical
factor belongs to the sum over indistinguishable final states.

\SupplementarySolution{18}{Multiplicity in a Sigma Decay}
Differentiating \(-a\sigma\pi_i\pi_i\) gives
\[
 i\mathcal M(\sigma\to\pi_i\pi_j)=-2ia\,\delta_{ij}.
\]
For one fixed flavor the daughters are identical.  The massless
two-body formula therefore gives
\[
 \Gamma_i=\frac1{2!}\frac{|2a|^2}{16\pi M}
 =\frac{a^2}{8\pi M}.
\]
Only equal-flavor pairs occur, but there are \(N-1\) possible flavors.
Hence
\[
 \Gamma_{\rm total}=\frac{(N-1)a^2}{8\pi M}.
\]
The factors \(2\), \(1/2!\), and \(N-1\) originate respectively from the
two identical fields at the vertex, identical-particle phase space, and
the sum over orthogonal flavor channels.

\SupplementarySolution{19}{A Connected \(3\to3\) Tree Amplitude}
The connected contribution first occurs in the second Dyson term:
\[
 \frac{(-i)^2}{2!}\left(\frac{\lambda}{4!}\right)^2
 \int d^4x\,d^4y\,
 T\{\phi^4(x)\phi^4(y)\}.
\]
Three external legs contract at each vertex and the remaining field at
each vertex forms the internal line.  Taking all six external momenta
\(k_a\) incoming, a chosen split \(A\cup A^c\), \(|A|=3\), contributes
\[
 i\mathcal M_A=(-i\lambda)^2
 \frac{i}{q_A^2-m^2+i0},\qquad
 q_A=\sum_{a\in A}k_a
\]
in the mostly-minus propagator convention.  The replacement
\(A\leftrightarrow A^c\) describes the same graph, so the number of
distinct labeled trees is
\[
 \frac12\binom63=10.
\]
The Dyson \(1/2!\) is cancelled by interchanging the two vertices, and
the \(4!\) factors are cancelled by their Wick contractions.

\SupplementarySolution{20}{Exponentiation of Vacuum Bubbles}
Let \(C_1\) be the sum of connected one-vertex vacuum bubbles, including
their integrals and symmetry factors, and \(C_2\) the connected
two-vertex contribution.  Wick's theorem organizes the vacuum amplitude
as
\[
 \braket{0}{S|0}
 =1+C_1+\left(C_2+\frac12C_1^2\right)+O(\lambda^3).
\]
The \(1/2\) divides permutations of the two identical disconnected
components.  This is precisely
\[
 \exp(C_1+C_2+\cdots)
 =1+C_1+C_2+\frac12C_1^2+\cdots.
\]
The same component-counting argument works at every order: a collection
of \(n_r\) connected bubbles of type \(r\) carries \(1/n_r!\).
Normalized matrix elements divide by \(\braket{0}{S|0}\), so the common
vacuum exponential cancels and no vacuum bubbles remain in scattering
amplitudes.

\SupplementarySolution{21}{An \(O(3)\)-Symmetric Scattering Tensor}
An invariant rank-four tensor with the required Bose symmetries is a
multiple of
\[
 \delta_{ij}\delta_{kl}+\delta_{ik}\delta_{jl}
 +\delta_{il}\delta_{jk}.
\]
Four differentiation of
\(-(\lambda/8)(\phi_a\phi_a)^2\) fixes the vertex to
\[
 -i\lambda\left(
 \delta_{ij}\delta_{kl}+\delta_{ik}\delta_{jl}
 +\delta_{il}\delta_{jk}\right).
\]
Thus
\[
 \mathcal M_{ij\to kl}
 =-\lambda\left(
 \delta_{ij}\delta_{kl}+\delta_{ik}\delta_{jl}
 +\delta_{il}\delta_{jk}\right).
\]
The tensor is unchanged by any \(O(3)\) rotation and is symmetric under
interchange within either the incoming or outgoing pair.

\SupplementarySolution{22}{The Exchange Minus Sign}
An ordered two-fermion out-state is created by
\(b^\dagger_3b^\dagger_4\ket0\).  Interchanging its labels gives
\[
 b^\dagger_4b^\dagger_3\ket0
 =-b^\dagger_3b^\dagger_4\ket0.
\]
Accordingly the two contractions that attach the outgoing fermions to
the incoming lines enter with opposite signs.  For scalar Yukawa
exchange one may write, up to a common coupling and phase convention,
\[
 \mathcal M=
 g^2\left[
 \frac{(\bar u_3u_1)(\bar u_4u_2)}{t-\mu^2}
 -\frac{(\bar u_4u_1)(\bar u_3u_2)}{u-\mu^2}
 \right].
\]
Under \(3\leftrightarrow4\), the two terms exchange and the amplitude
changes sign.  This antisymmetry, rather than a manually appended rule,
is the direct consequence of the external fermionic Fock states.

\SupplementarySolution{23}{A Step Function and the \(i0\) Prescription}
Put \(x=\omega-\omega_0\).  For \(\epsilon>0\),
\[
 I_\epsilon(x)=\int_0^\infty dt\,e^{ixt-\epsilon t}
 =\frac1{\epsilon-ix}=\frac{i}{x+i\epsilon}.
\]
The Sokhotski--Plemelj formula gives
\[
 \frac1{x+i0}=\operatorname{P.V.}\frac1x-i\pi\delta(x),
\]
and hence
\[
 \lim_{\epsilon\to0^+}I_\epsilon(x)
 =i\,\operatorname{P.V.}\frac1x+\pi\delta(x).
\]
The delta part comes from the indefinitely long coherent time interval;
the principal-value part is the dispersive contribution.

\SupplementarySolution{24}{Ordered Times from Energy Denominators}
For intermediate energies \(E_a,E_b\),
\[
 \frac1{E_i-E_a+i0}\frac1{E_i-E_b+i0}
 =(-i)^2\int_0^\infty d\tau_1d\tau_2\,
 e^{i(E_i-E_a)\tau_1+i(E_i-E_b)\tau_2}.
\]
After inserting the overall energy delta function, choose an unrestricted
latest time \(t_1\) and set
\[
 t_2=t_1-\tau_1,\qquad
 t_3=t_2-\tau_2.
\]
The Jacobian is one, and \(\tau_1,\tau_2\geq0\) is equivalent to
\[
 -\infty<t_3\leq t_2\leq t_1<\infty.
\]
Inserting complete energy eigenstates between three interactions makes
the exponentials telescope into the interaction-picture phases of
\(H_I(t_1)H_I(t_2)H_I(t_3)\).  The two denominators therefore become
the three nested time integrals in the third-order Dyson term.

\SupplementarySolution{25}{Elastic Partial-Wave Unitarity}
Orthogonality of angular momentum channels makes unitarity diagonal:
\(|S_\ell|=1\).  With
\[
 S_\ell=1+2ikf_\ell,
\]
the identity \(S_\ell^*S_\ell=1\) gives
\[
 \operatorname{Im}f_\ell=k|f_\ell|^2.
\]
Every unit-modulus complex number can be written
\(S_\ell=e^{2i\delta_\ell}\) with real \(\delta_\ell\).  Hence
\[
 f_\ell=\frac{e^{2i\delta_\ell}-1}{2ik}
 =\frac{e^{i\delta_\ell}\sin\delta_\ell}{k}.
\]
This form satisfies the unitarity relation identically and makes the
phase-shift ambiguity \(\delta_\ell\sim\delta_\ell+\pi\) explicit.

\SupplementarySolution{26}{Partial-Wave Cross Sections and the Unitarity Bound}
Using
\(\int d\Omega\,P_\ell(\cos\theta)P_{\ell'}(\cos\theta)
=4\pi\delta_{\ell\ell'}/(2\ell+1)\), one finds
\[
 \sigma_{\rm el}
 =4\pi\sum_\ell(2\ell+1)|f_\ell|^2
 =\frac{4\pi}{k^2}
 \sum_\ell(2\ell+1)\sin^2\delta_\ell.
\]
For purely elastic scattering this also equals the total cross section,
as follows independently from the forward optical theorem.  Since
\(\sin^2\delta_\ell\leq1\), one partial wave obeys
\[
 \sigma_\ell\leq\frac{4\pi}{k^2}(2\ell+1).
\]
The bound is saturated at
\(\delta_\ell=\pi/2\pmod{\pi}\).

\SupplementarySolution{27}{The One-Channel \(K\)-Matrix}
Equating the two representations,
\[
 e^{2i\delta}=\frac{1-iK}{1+iK},
\]
and solving gives
\[
 K=-\tan\delta.
\]
The minus sign follows from the chapter's convention
\(S=1-2iK+\cdots\); the opposite convention for the transition operator
would give \(K=\tan\delta\).  At a resonance the phase passes through
\(\delta=\pi/2\), so \(K\) has a real-axis pole while \(S=-1\) remains
finite.  If \(K=K^*\), then
\[
 S^*=\frac{1+iK}{1-iK}=S^{-1},
\]
which proves one-channel unitarity without further constraints.

\SupplementarySolution{28}{Inelasticity in a Partial Wave}
Unitarity of the full multichannel matrix implies that the squared norm
of any one diagonal entry cannot exceed one.  Thus
\[
 S_\ell=\eta_\ell e^{2i\delta_\ell},\qquad
 0\leq\eta_\ell\leq1.
\]
With \(f_\ell=(S_\ell-1)/(2ik)\), the three standard partial-wave
cross sections are
\[
 \begin{split}
 \sigma_{{\rm el},\ell}
 &=\frac{\pi}{k^2}(2\ell+1)
 |1-\eta_\ell e^{2i\delta_\ell}|^2,\\
 \sigma_{{\rm reac},\ell}
 &=\frac{\pi}{k^2}(2\ell+1)(1-\eta_\ell^2),\\
 \sigma_{{\rm tot},\ell}
 &=\frac{2\pi}{k^2}(2\ell+1)
 (1-\eta_\ell\cos2\delta_\ell).
 \end{split}
\]
Direct expansion verifies
\(\sigma_{\rm tot}=\sigma_{\rm el}+\sigma_{\rm reac}\).

\SupplementarySolution{29}{Breit--Wigner Phase Motion}
For real \(E\), numerator and denominator are complex conjugates, so
\(|S_\ell|=1\).  Put \(x=E-E_R\) and \(\gamma=\Gamma/2\).  A continuous
choice increasing from \(0\) to \(\pi\) obeys
\[
 \tan\delta_\ell=-\frac{\gamma}{x},
\qquad
 \delta_\ell(E_R)=\frac{\pi}{2}.
\]
Therefore
\[
 \sin^2\delta_\ell
 =\frac{\gamma^2}{x^2+\gamma^2}
 =\frac{\Gamma^2/4}
 {(E-E_R)^2+\Gamma^2/4}.
\]
Using
\(\lim_{\gamma\to0^+}\gamma/(x^2+\gamma^2)=\pi\delta(x)\),
\[
 \sin^2\delta_\ell
 \longrightarrow\pi\gamma\,\delta(E-E_R)
 =\frac{\pi\Gamma}{2}\delta(E-E_R).
\]
The peak reaches the elastic unitarity limit independently of its width,
while its integrated area is proportional to \(\Gamma\).

\SupplementarySolution{30}{Charge Channels from Isospin}
The state \(\ket{\pi^+p}\) has \(I_3=3/2\), hence
\[
 A(\pi^+p\to\pi^+p)=A_{3/2}.
\]
For \(I_3=-1/2\), a convenient Clebsch--Gordan convention is
\[
 \begin{split}
 \ket{\pi^-p}
 &=\sqrt{\frac13}\ket{\tfrac32,-\tfrac12}
 -\sqrt{\frac23}\ket{\tfrac12,-\tfrac12},\\
 \ket{\pi^0n}
 &=\sqrt{\frac23}\ket{\tfrac32,-\tfrac12}
 +\sqrt{\frac13}\ket{\tfrac12,-\tfrac12}.
 \end{split}
\]
An isospin-invariant amplitude is diagonal in \(I\), so
\[
 \begin{split}
 A(\pi^-p\to\pi^-p)
 &=\frac13A_{3/2}+\frac23A_{1/2},\\
 A(\pi^-p\to\pi^0n)
 &=\frac{\sqrt2}{3}(A_{3/2}-A_{1/2}).
 \end{split}
\]
The sign of the charge-exchange line changes if the phase of
\(\ket{\pi^0n}\) is reversed, with no observable consequence.  Finally,
\[
 |A_{\rm el}|^2+|A_{\rm cx}|^2
 =\frac13|A_{3/2}|^2+\frac23|A_{1/2}|^2,
\]
because the interference terms cancel.  This is precisely preservation
of norm by the Clebsch--Gordan change of basis.
\end{enumerate}
